\documentclass[11pt,a4paper]{article}
\usepackage[T1]{fontenc}
\usepackage[utf8]{inputenc}
\usepackage[polish]{babel}
\usepackage{amsmath,amssymb,booktabs,geometry}
\geometry{margin=2.5cm}
\newcommand{\expv}[1][]{\langle O\rangle}
\title{Twirling w formalizmie macierzy procesu $\chi$\\
\large materiał uzupełniający do raportu QEM}
\author{Karina Leśkiewicz \\ \small KN AXION, Politechnika Wrocławska}
\date{\today}
\begin{document}
\maketitle
\section*{Wyprowadzenie działania twirlingu w formalizmie $\chi$}

Dowolny kanał kwantowy $\mathcal{E}$ działający na $n$ kubitach można zapisać w~bazie operatorów
Pauliego $\{P_\alpha\}$ jako
\begin{equation}
  \mathcal{E}(\rho) \;=\; \sum_{\alpha\beta} \chi_{\alpha\beta}\, P_\alpha \rho\, P_\beta^{\dagger},
  \label{eq:chi-def}
\end{equation}
gdzie $\chi$ (macierz procesu) jest hermitowska i~dodatnio półokreślona. Interpretacja jej elementów
jest kluczowa:
\begin{itemize}
  \item elementy \textbf{diagonalne} $\chi_{\alpha\alpha}$ to prawdopodobieństwa, z~jakimi kanał
        stosuje błąd Pauliego $P_\alpha$ --- składowa \emph{stochastyczna};
  \item elementy \textbf{pozadiagonalne} $\chi_{\alpha\beta}$, $\alpha\neq\beta$, opisują
        \emph{interferencję} między różnymi ścieżkami błędu --- składową \emph{koherentną},
        odpowiedzialną za kumulację amplitudową.
\end{itemize}
Błąd czysto unitarny (np.\ systematyczne przekręcenie $U_\theta = e^{-i\theta P/2}$) ma znaczące
elementy pozadiagonalne; kanał depolaryzujący ma wyłącznie diagonalę.

\paragraph{Twirl Pauliego.} Otoczmy kanał losową bramką Pauliego $P_k$ wybraną jednostajnie
z~$n$-kubitowej grupy Pauliego i~uśrednijmy po tym wyborze:
\begin{equation}
  \mathcal{E}_{\text{tw}}(\rho) \;=\; \frac{1}{4^n}\sum_k
  P_k^{\dagger}\,\mathcal{E}\bigl(P_k \rho P_k^{\dagger}\bigr)\,P_k .
\end{equation}
Podstawiając~\eqref{eq:chi-def} i~korzystając z~tego, że operatory Pauliego albo komutują, albo
antykomutują ($P_k P_\alpha = (-1)^{\langle k,\alpha\rangle} P_\alpha P_k$ dla pewnej funkcji
symplektycznej $\langle\cdot,\cdot\rangle$ o~wartościach w~$\{0,1\}$), otrzymujemy
\begin{equation}
  \mathcal{E}_{\text{tw}}(\rho) \;=\; \sum_{\alpha\beta}
  \underbrace{\Bigl[\frac{1}{4^n}\sum_k (-1)^{\langle k,\alpha\rangle + \langle k,\beta\rangle}\Bigr]}
  _{\textstyle =\;\delta_{\alpha\beta}}
  \chi_{\alpha\beta}\, P_\alpha \rho\, P_\beta^{\dagger} .
\end{equation}
Czynnik w~nawiasie jest średnią znaków po całej grupie Pauliego. Dla $\alpha=\beta$ wykładnik jest
parzysty i~wszystkie wyrazy równe $+1$, więc średnia wynosi 1. Dla $\alpha\neq\beta$ istnieje
element grupy komutujący z~jednym operatorem, a~antykomutujący z~drugim, wobec czego znaki
$+1$ i~$-1$ występują dokładnie po równo i~średnia znika (ortogonalność charakterów). W~rezultacie
\begin{equation}
  \boxed{\;\mathcal{E}_{\text{tw}}(\rho) \;=\; \sum_{\alpha} \chi_{\alpha\alpha}\,
  P_\alpha \rho\, P_\alpha^{\dagger}\;}
\end{equation}
--- kanał po twirlingu jest \emph{kanałem Pauliego}: zachowuje diagonalę $\chi$ (a~więc i~całkowitą
wierność procesu), lecz \textbf{zeruje wszystkie elementy pozadiagonalne}. Błąd koherentny zostaje
zamieniony na stochastyczny o~tej samej sile.

\paragraph{Realizacja praktyczna (randomized compiling).} Bezpośrednie wykonanie powyższego
uśrednienia wymagałoby dodatkowych bramek. W~randomized compiling~\cite{wallman2016,hashim2021}
wykorzystuje się fakt, że bramki Cliffordowskie przekształcają operatory Pauliego w~operatory
Pauliego: bramkę dwukubitową $G$ otacza się losową bramką $A$ przed nią i~\emph{kompensatorem}
\begin{equation}
  B \;=\; G\,A\,G^{\dagger},
\end{equation}
tak że $B\,G\,A = G$ z~dokładnością do fazy globalnej --- działanie idealne pozostaje niezmienione,
a~uśrednienie po $M$ losowaniach realizuje twirl. Ponieważ $A$ i~$B$ są pojedynczymi bramkami
jednokubitowymi, narzut jest niewielki: w~naszych obwodach wynosi $8$--$13\%$ dodatkowych bramek
jednokubitowych przy zmianie głębokości nieprzekraczającej $5\%$.

% =====================================================================================

\end{document}
